\section{Propuesta de Política de Seguridad de la información}
La empresa deberá implementar una Política de Seguridad de la Información enfocada en proteger la confidencialidad, integridad y disponibilidad de los datos relacionados con los procesos de pesaje, facturación, pagos y control administrativo. Esta política aplicará para todos los empleados que utilicen equipos informáticos, sistemas de red y dispositivos de almacenamiento.

\subsection*{1. Control de acceso y uso de cuentas}
\begin{itemize}
    \item Cada empleado deberá contar con un usuario y contraseña únicos para acceder a los sistemas de operación y administración.
    \item Queda prohibido compartir contraseñas entre empleados.
    \item Las contraseñas deberán cambiarse periódicamente y contar con una complejidad mínima.
    \item Se deberán activar registros de auditoría (logs) para detectar modificaciones.
\end{itemize}

\subsection*{2. Principios de diseño seguro}
\begin{itemize}
    \item Los sistemas utilizados por la empresa deberán diseñarse considerando la seguridad desde el inicio y no como una característica adicional.
    \item Los equipos críticos deberán estar protegidos físicamente contra polvo, calor y variaciones eléctricas.
    \item Se implementarán mecanismos de respaldo y recuperación de información para asegurar la continuidad operativa.
    \item Los sistemas deberán operar bajo el principio de mínimo privilegio, permitiendo únicamente los accesos necesarios para cada usuario.
    \item Los programas y sistemas utilizados deberán mantenerse actualizados para reducir vulnerabilidades.
\end{itemize}

\subsection*{3. Programación en defensa}
\begin{itemize}
    \item Todo software utilizado o desarrollado para la empresa deberá validar entradas de datos para evitar errores o manipulaciones.
    \item Se deberán implementar controles que detecten modificaciones indebidas en bases de datos o registros de pesaje.
    \item El sistema deberá generar alertas ante intentos de acceso no autorizados.
    \item Se recomienda utilizar antivirus, firewall y monitoreo constante para prevenir malware y ransomware.
    \item Los desarrolladores o técnicos responsables deberán aplicar prácticas seguras de codificación y pruebas antes de implementar cambios.
\end{itemize}

\subsection*{4. Amenazas y ataques para redes}
\begin{itemize}
    \item Se reconoce como amenazas principales: ransomware, phishing, robo de información, malware y accesos no autorizados.
    \item La empresa deberá restringir el uso de dispositivos USB externos no autorizados.
    \item Se implementará antivirus actualizado y monitoreo de red para detectar actividades sospechosas.
    \item Los empleados recibirán capacitación básica para identificar correos fraudulentos y ataques de ingeniería social.
    \item Se deberán utilizar redes seguras y segmentadas para separar equipos administrativos y operativos.
\end{itemize}

\subsection*{5. Integridad y confidencialidad de los mensajes en la red}
\begin{itemize}
    \item Toda información relacionada con proveedores, facturación y datos fiscales deberá protegerse mediante mecanismos de cifrado.
    \item Los archivos sensibles no deberán enviarse mediante medios inseguros o públicos.
    \item Se deberán realizar respaldos periódicos de la información en medios seguros y protegidos.
    \item Los datos transmitidos dentro de la red deberán mantenerse íntegros para evitar modificaciones no autorizadas.
\end{itemize}

\subsection*{6. Protocolos seguros de comunicación}
\begin{itemize}
    \item La empresa deberá utilizar protocolos seguros como HTTPS, SFTP o VPN para la transferencia de información sensible.
    \item Se evitará el uso de software obsoleto o sin soporte técnico.
    \item Los equipos conectados a internet deberán contar con configuraciones seguras y actualizaciones constantes.
    \item El acceso remoto a sistemas administrativos solo deberá realizarse mediante conexiones seguras y autenticadas.
\end{itemize}

\subsection*{7. Continuidad operativa y respaldo}
\begin{itemize}
    \item Se deberán implementar No-Breaks (UPS) para proteger los equipos críticos ante fallas eléctricas.
    \item Se realizarán respaldos periódicos en nube y dispositivos externos protegidos.
    \item La empresa deberá contar con procedimientos de recuperación ante fallas o pérdida de información.
\end{itemize}
